\section*{Capítulo \ref{ch:b2:orbital-forcing} --- O forçamento orbital e as eras glaciais}

\begin{solution}{exo:b2:orbital-forcing:1}
(a) $1.0167/0.9833 = 1.0340$: o afélio está \qty{3.4}{\percent} mais longe que o periélio.
(b) $a(1 + e) - a(1 - e) = 2ea = 0.0334\,a$, \qty{3.34}{\percent} de $a$, perto dos 3.4 por cento
habitualmente citados. (c) A luz solar diminui com o inverso do quadrado da distância: $1.0340^2 = 1.069$,
cerca de \qty{7}{\percent} mais luz solar no periélio.
\end{solution}

\begin{solution}{exo:b2:orbital-forcing:2}
(a) Hoje $0.017236 \times \sin\ang{101.37} = 0.017236 \times 0.9804 = 0.01690$; há \num{10000} anos
$0.019971 \times \sin\ang{293.86} = 0.019971 \times (-0.9145) = -0.01826$; há \num{6000} anos
$0.019249 \times \sin\ang{359.99} \approx 0$, os valores das tabelas. (b) A longitude do Sol no
periélio é $\varpi + \ang{180}$: \ang{281.37} hoje, \ang{473.86} $-$ \ang{360} $=$ \ang{113.86}
há \num{10000} anos, \ang{179.99} há \num{6000} anos. (c) Há \num{10000} anos o periélio
caía \ang{24} depois do solstício de junho, o mais próximo dele dos três, no verão boreal;
há \num{6000} anos caía no equinócio de setembro, onde a \omterm{def:b2:orbital-forcing:elements}{precessão climática} é nula.
\end{solution}

\begin{solution}{exo:b2:orbital-forcing:3}
(a) As frequências somam-se: $1/25.77 + 1/112 = 0.03880 + 0.00893 = \qty{0.04773}{kyr^{-1}}$, um período
de $1/0.04773 = 21.0$ mil anos. (b) O resultado fica dentro do intervalo dos períodos principais,
19.0 a 23.7 mil anos. Um só número não pode reproduzi-los: a rotação da órbita
é ela mesma a soma de vários movimentos de frequências diferentes, cada um dos quais se combina com a
precessão do eixo para dar um período da \omterm{def:b2:orbital-forcing:elements}{precessão climática}; a média é os
\num{23000} anos habitualmente citados.
\end{solution}

\begin{solution}{exo:b2:orbital-forcing:4}
$\sqrt{1 - 0.0167^2} = 0.99986$. Para \ang{22.1}: $1360 \times 0.37622/(\pi \times 0.99986) =
\qty{162.9}{W.m^{-2}}$; para \ang{24.5}: $1360 \times 0.41469/(\pi \times 0.99986) =
\qty{179.5}{W.m^{-2}}$. A razão é $\sin\ang{24.5}/\sin\ang{22.1} = 1.102$: a média anual no
polo cresce \qty{10.2}{\percent}, cerca de \qty{16.7}{W.m^{-2}}, entre os dois extremos.
\end{solution}

\begin{solution}{exo:b2:orbital-forcing:5}
(a) No solstício de junho, tem-se $(a/r)^2 = \bigl((1 - e\sin\varpi)/(1 - e^2)\bigr)^2$. Hoje, $e\sin\varpi
= 0.016898$ e $(0.983102/0.999703)^2 = 0.9671$, contra $1 - 0.033796 = 0.9662$;
há \num{116000} anos $e\sin\varpi = 0.043939$ e $(0.956061/0.998065)^2 = 0.9176$, contra $1 -
0.087878 = 0.9121$. (b) No solstício de dezembro $(a/r)^2 = \bigl((1 + e\sin\varpi)/(1 -
e^2)\bigr)^2$: hoje $1.0347$ contra $1.0338$, há \num{116000} anos $1.0940$ contra $1.0879$.
(c) Os termos desprezados são de ordem $e^2$: $(1 - x)^2/(1 - e^2)^2 \approx 1 - 2x + x^2 + 2e^2$
com $x = e\sin\varpi$, aqui cerca de $3e^2$. Valem $3 \times 0.0172^2 = 0.0009$ hoje e $3
\times 0.0440^2 = 0.0058$ há \num{116000} anos, perto das diferenças encontradas, $0.0009$ e
$0.0055$: a aproximação é seis vezes pior quando a excentricidade é duas vezes e meia
maior.
\end{solution}

\begin{solution}{exo:b2:orbital-forcing:6}
No solstício de junho $\delta = \varepsilon = \ang{23.446}$. $\cos h_0 = -\tan\ang{65}\tan\ang{23.446}
= -2.14451 \times 0.43369 = -0.93006$, logo $h_0 = \ang{158.44} = \qty{2.7654}{rad}$. O dia dura
$2 \times 158.44/15 = 21.13$~horas (a Terra gira \ang{15} por hora). Com $\sin\ang{65} = 0.90631$,
$\cos\ang{65} = 0.42262$, $\sin\delta = 0.39788$, $\cos\delta = 0.91744$ e $\sin h_0 = 0.36742$,
o parêntese vale $2.7654 \times 0.90631 \times 0.39788 + 0.42262 \times 0.91744 \times 0.36742 = 0.99720
+ 0.14246 = 1.13966$. À distância média $\bar Q = (1360/\pi) \times 1.13966 = 432.90 \times 1.13966
= \qty{493.4}{W.m^{-2}}$; com o fator de distância do \cref{exo:b2:orbital-forcing:5}, $493.4 \times
0.9671 = \qty{477.1}{W.m^{-2}}$. Como verificação, as mesmas etapas a \ang{60} N dão \qty{475.96}{W.m^{-2}}, onde as tabelas dão \num{475.95}.
\end{solution}

\begin{solution}{exo:b2:orbital-forcing:7}
(a) O equinócio de março é $\lambda = 0$, o equinócio de setembro $\lambda = \pi$. Com $\sqrt{1 - e^2}
\approx 1$, a duração é $\frac{T}{2\pi}\int_0^{\pi}\bigl(1 + 2e\cos(\lambda -
\varpi)\bigr)\,\mathrm d\lambda = \frac{T}{2\pi}\bigl(\pi + 2e[\sin(\lambda - \varpi)]_0^{\pi}\bigr)
= \frac{T}{2\pi}(\pi + 4e\sin\varpi) = \frac T2 + \frac{2T}{\pi}e\sin\varpi$. (b) O outro semestre
é mais curto na mesma quantidade, de modo que a diferença é $(4T/\pi)\,e\sin\varpi = (4 \times
365.2422/\pi) \times 0.016898 = 465.04 \times 0.016898 = 7.86$~dias. (c) Primavera e verão juntos
excedem outono e inverno em $4.5 + 3 = 7.5$~dias; as duas diferenças citadas são arredondadas, e a
concordância fica dentro desse arredondamento.
\end{solution}

\begin{solution}{exo:b2:orbital-forcing:8}
(a) Para $e = 0.0046$: $\sqrt{1 - e^2} = 0.999989$ e $340/0.999989 = \qty{340.004}{W.m^{-2}}$. Para
$e = 0.0565$: $\sqrt{1 - 0.003192} = 0.998403$ e $340/0.998403 = \qty{340.544}{W.m^{-2}}$. (b) A
diferença é \qty{0.54}{W.m^{-2}}, \qty{0.16}{\percent}. (c) A insolação de meados de julho a \ang{65} N
oscilou $497 - 388 = \qty{109}{W.m^{-2}}$, cerca de \qty{25}{\percent} de sua média e duzentas
vezes a variação da média global anual: a órbita age sobre o clima deslocando a luz solar entre
estações e latitudes, não alterando o total.
\end{solution}

\begin{solution}{exo:b2:orbital-forcing:9}
Com $\omega = 2\pi/P$ e $\tau = \qty{10}{kyr}$: para $P = 23$, $\omega\tau = 2.732$, uma amplitude
relativa $1/\sqrt{1 + 7.464} = 0.344$ e um atraso $\arctan(2.732)/\omega = 1.2199/0.2732 = 4.47$~kyr;
para $P = 41$, $\omega\tau = 1.532$, amplitude $0.546$, atraso $0.9926/0.1532 = 6.48$~kyr; para $P =
100$, $\omega\tau = 0.628$, amplitude $0.847$, atraso $0.5610/0.0628 = 8.93$~kyr. O manto de gelo
favorece o forçamento lento: segue \qty{85}{\percent} de um ciclo de \num{100000} anos, mas só um
terço de um ciclo de precessão, pois tem pouco tempo para crescer ou diminuir antes que um forçamento rápido
se inverta. Ele age como um filtro passa-baixa.
\end{solution}

\begin{solution}{exo:b2:orbital-forcing:10}
(a) A equação é linear: se $V_1$ responde a $A_1\sin\omega_1t$ e $V_2$ responde a
$A_2\sin\omega_2t$, $V_1 + V_2$ responde à sua soma. Pela \cref{prop:b2:orbital-forcing:lag} $V_1$ e
$V_2$ são senoides de frequências $\omega_1$ e $\omega_2$ (uma vez extinto o transitório), de modo que
a resposta não contém nenhuma outra frequência. (b) Com $V \approx B_1\sin(\omega_1t - \phi_1) +
B_2\sin(\omega_2t - \phi_2)$,
\[
\begin{aligned}
  V^2 &= \tfrac12(B_1^2 + B_2^2) - \tfrac12B_1^2\cos(2\omega_1t - 2\phi_1) -
  \tfrac12B_2^2\cos(2\omega_2t - 2\phi_2)\\
  &\quad + B_1B_2\bigl[\cos\bigl((\omega_1 - \omega_2)t - \phi_1 + \phi_2\bigr) -
  \cos\bigl((\omega_1 + \omega_2)t - \phi_1 - \phi_2\bigr)\bigr],
\end{aligned}
\]
de modo que o termo $kV^2$ força a equação nas frequências diferença e soma, e nas frequências
duplas. (c) $1/(1/19.0 - 1/23.7) = 1/(0.05263 - 0.04219) = 95.8$ mil anos; o par 19.0
e 22.4 dá 125 mil anos. As frequências diferença dos termos de precessão caem na banda da
excentricidade, como devem, já que a \omterm{def:b2:orbital-forcing:elements}{precessão climática} é $e\sin\varpi$. Um manto de gelo
não linear poderia portanto extrair da precessão um ritmo de \num{100000} anos; um linear
não. Os ciclos de \num{100000} anos são indício de um comportamento não linear, embora este cálculo
sozinho não mostre qual.
\end{solution}

\begin{solution}{exo:b2:orbital-forcing:11}
(a) Como na demonstração da \cref{prop:b2:orbital-forcing:annual}, $(a/r)^2\,\mathrm dt =
T\,\mathrm d\lambda/(2\pi\sqrt{1 - e^2})$, e $\sin\delta = \sin\varepsilon\sin\lambda$ é positivo
para $0 < \lambda < \pi$. Logo $E = \frac{S_0T\sin\varepsilon}{2\pi\sqrt{1 - e^2}}\int_0^{\pi}\sin\lambda
\,\mathrm d\lambda = S_0T\sin\varepsilon/(\pi\sqrt{1 - e^2})$. (b) $T = 365.25 \times 86400 =
\qty{3.1558e7}{s}$, $S_0T/\pi = \qty{1.3661e10}{J.m^{-2}}$; para \ang{22.1}, $E = 1.3661 \times 10^{10}
\times 0.37622/0.99986 = \qty{5.14}{GJ.m^{-2}}$; para \ang{24.5}, $\qty{5.67}{GJ.m^{-2}}$,
\qty{10}{\percent} a mais. (c) A precessão fixa como a mesma energia se distribui no tempo: um verão
no periélio é mais curto e mais intenso, um verão no afélio mais longo e mais fraco. O derretimento só começa
acima de \qty{0}{\celsius}, de modo que um verão curto e intenso pode derreter mais que um longo e fresco
que receba o mesmo total.
\end{solution}

\begin{solution}{exo:b2:orbital-forcing:12}
(a) $3.93\log_2(279.4/189.7) = 3.93\log_2 1.4729 = 3.93 \times 0.5586 = \qty{2.20}{W.m^{-2}}$. (b) A
menor subida é de 194 para \qty{266}{ppm}: $3.93\log_2 1.3711 = 3.93 \times 0.4554 =
\qty{1.79}{W.m^{-2}}$; a maior de 188 para \qty{282}{ppm}: $3.93\log_2 1.5 = 3.93 \times 0.5850 =
\qty{2.30}{W.m^{-2}}$. (c) O forçamento do dióxido de carbono é quatro vezes a maior variação da
insolação global anual, \qty{0.54}{W.m^{-2}}, e age em toda parte e o ano todo; a insolação de verão
a \ang{65} N oscila cerca de \qty{109}{W.m^{-2}}, mas numa latitude e numa estação. A
órbita é o gatilho, através dos verões dos mantos de gelo do norte; o dióxido de carbono, com a
\omterm{def:b2:precambrian-earth:ice-albedo}{retroalimentação gelo--albedo}, é um amplificador que espalha a mudança sobre o globo.
\end{solution}

\begin{solution}{pb:b2:orbital-forcing:1}
\textbf{1.} Então $0.0440 \times \sin\ang{92.7} = 0.0440 \times 0.99889 = 0.04395$; agora $0.0172 \times
\sin\ang{101.4} = 0.0172 \times 0.98027 = 0.01686$, cerca de 2.6 vezes menos.
\textbf{2.} Então $1 - 0.0440 = 0.956$ e $1 + 0.0440 = 1.044$; agora $0.9828$ e $1.0172$.
\textbf{3.} Em $\lambda = \varpi + \ang{180}$: então \ang{272.7}, agora \ang{281.4}, ambos perto do
solstício de dezembro ($\lambda = \ang{270}$); então a Terra passava pelo periélio menos de três dias
depois dele, e pelo afélio no solstício de junho.
\textbf{4.} Então $(1.044/0.956)^2 = 1.0921^2 = 1.193$, agora $(1.0172/0.9828)^2 = 1.0350^2 = 1.071$:
\qty{19}{\percent} mais luz solar no periélio então, \qty{7}{\percent} agora.
\textbf{5.} Menor, \ang{22.52} contra \ang{23.45}: o sol de verão ficava mais baixo nas altas latitudes,
e as altas latitudes recebiam menos luz solar no verão e ao longo do ano.
\textbf{6.} Então $340/\sqrt{1 - 0.001936} = 340/0.999032 = \qty{340.33}{W.m^{-2}}$; agora $340/0.999852
= \qty{340.05}{W.m^{-2}}$. A diferença, \qty{0.28}{W.m^{-2}} ou menos de uma parte por mil,
é insignificante.
\textbf{7.} Em $\lambda = \ang{90}$, $\cos(\ang{90} - \varpi) = \sin\varpi$ na
\cref{prop:b2:orbital-forcing:distance}.
\textbf{8.} Então $(1 - 0.001936)/(1 - 0.04395) = 0.998064/0.956049 = 1.04395$; agora $(1 -
0.000296)/(1 - 0.01686) = 0.999704/0.983139 = 1.01685$.
\textbf{9.} Então $1/1.04395^2 = 0.91758$; agora $1/1.01685^2 = 0.96714$.
\textbf{10.} Então $1 - 2 \times 0.04395 = 0.91210$; agora $1 - 2 \times 0.01686 = 0.96628$: a
aproximação erra em $0.0055$ então, $0.0009$ agora.
\textbf{11.} Então $1 - 0.91758 = \qty{8.2}{\percent}$ mais fraca, agora \qty{3.3}{\percent}.
\textbf{12.} No solstício $\delta = \varepsilon$: \ang{22.52} então, \ang{23.45} agora.
\textbf{13.} Então $\cos h_0 = -2.14451 \times 0.41462 = -0.88916$, $h_0 = \ang{152.77} =
\qty{2.6663}{rad}$; agora $\cos h_0 = -2.14451 \times 0.43378 = -0.93023$, $h_0 = \ang{158.47} =
\qty{2.7659}{rad}$.
\textbf{14.} $2h_0/\ang{15}$ por hora: então $20.37$~horas, agora $21.13$~horas.
\textbf{15.} Então $2.6663 \times 0.90631 \times 0.38301 + 0.42262 \times 0.92375 \times 0.45759 =
0.92554 + 0.17864 = 1.10418$ (\num{1.10417} sem arredondar), e $432.90 \times 1.10417 =
\qty{478.0}{W.m^{-2}}$; agora $2.7659 \times 0.90631 \times 0.39795 + 0.42262 \times 0.91741 \times
0.36697 = 0.99757 + 0.14228 = 1.13985$ (\num{1.13982} sem arredondar), e $432.90 \times 1.13982 = \qty{493.4}{W.m^{-2}}$.
\textbf{16.} Então $478.0 \times 0.91758 = \qty{438.6}{W.m^{-2}}$; agora $493.4 \times 0.96714 =
\qty{477.2}{W.m^{-2}}$, \qty{38.6}{W.m^{-2}} a mais que então.
\textbf{17.} Em meados de julho a longitude do Sol é \ang{120}: $\sin\delta = \sin\varepsilon\sin\ang{120}$,
uma declinação de cerca de \ang{20} em vez de \ang{23}, logo um dia mais curto e um sol mais baixo: a
\ang{65} N a insolação média diária é máxima perto do solstício, quando a declinação é máxima.
\textbf{18.} Com a distância de hoje, a variação da obliquidade sozinha custa $(493.4 - 478.0) \times
0.96714 = \qty{14.9}{W.m^{-2}}$; com a obliquidade antiga, a distância sozinha custa $478.0 \times
(0.96714 - 0.91758) = \qty{23.7}{W.m^{-2}}$. As duas somam \qty{38.6}{W.m^{-2}}; a distância
pesou cerca de três quintos.
\textbf{19.} Depois de uma variação súbita o manto de gelo realizou $1 - \mathrm e^{-t/\tau}$ dela:
$1 - \mathrm e^{-0.5} = 0.39$ depois de \qty{5}{kyr}, $1 - \mathrm e^{-1} = 0.63$ depois de \qty{10}{kyr},
$1 - \mathrm e^{-2} = 0.86$ depois de \qty{20}{kyr}.
\textbf{20.} $\omega\tau = 2\pi \times 10/23 = 2.732$: o manto de gelo segue $1/\sqrt{1 + 2.732^2} =
0.34$ do forçamento, com $\arctan(2.732)/\omega = 4.5$~kyr de atraso.
\textbf{21.} Meio ciclo de precessão, uns 11 a \qty{12}{kyr}, é cerca de um tempo de resposta: antes que
o manto de gelo tenha percorrido dois terços do caminho, a insolação inverte-se e puxa-o de volta.
\textbf{22.} No modelo linear o gelo cresce e diminui ao mesmo ritmo para forçamentos do mesmo
tamanho: um forçamento que sobe e um que desce dão respostas simétricas. Uma acumulação lenta seguida de um colapso
rápido exige uma resposta que dependa do estado ou do sentido da variação, isto é, uma
resposta não linear.
\textbf{23.} O dióxido de carbono subiu $3.93\log_2(266/194) = 1.8$ a $3.93\log_2(282/188) =
\qty{2.3}{W.m^{-2}}$ do glacial ao \omterm{def:b2:orbital-forcing:glacials}{interglacial}: um forçamento sobre o globo inteiro e o ano
inteiro. Liberado quando a circulação atlântica enfraquece, ele leva a mudança dos verões do norte
ao planeta inteiro e, com o albedo do gelo que encolhe, amplifica-a.
\textbf{24.} Quando o último \omterm{def:b2:orbital-forcing:glacials}{interglacial} terminou, a insolação média diária do solstício de junho a
\ang{65} N era de cerca de \textbf{\qty{439}{W.m^{-2}}}, \qty{39}{W.m^{-2}} a menos que hoje e seu valor
mais baixo do ciclo.
\end{solution}
